\section{The two-band cotangent certificate}\label{sec:tangents}
For a matrix $M\in\SL_r$ let $e_j(M)$ be the $j$th elementary symmetric
function of its eigenvalues, so
$\det(zI-M)=\sum_{j=0}^r(-1)^je_j(M)z^{r-j}$. We use $1\leq j<r$;
$e_r=1$ is constant. For a group word $w$, these are regular invariant
functions on the representation variety and descend to its character
variety.

The use of centralizers of twisted products and their powers has a close
precedent in Pickrell--Xia \cite[Proposition 2.1.24, Lemma 2.1.34]{PX}.
Here we need the following explicit calculation at a family whose first
two handle matrices commute; no generic compact-group spanning theorem
is substituted for the calculation.

\begin{proposition}\label{prop:core}
Let $(A,B,Q)$ satisfy Lemma~\ref{lem:all-open}, and set $C=QC_\mu Q^{-1}$.
At the genus-two representation $(A,B,C,I)$, the common tangent kernel of
the characteristic coefficients of $A,C$ and of all the simple curves
$u_0(k,l),u_1(k,l)$, $k,l\geq1$, consists exactly of infinitesimal
simultaneous conjugations. The resulting character is a smooth
irreducible point, and these differentials span its cotangent space.
\end{proposition}
\begin{proof}
Since $A,C$ are regular semisimple, stationarity of their characteristic
coefficients says that their variations are tangent to their conjugacy
classes. Subtract an infinitesimal simultaneous conjugation to write a
common-kernel tangent vector as
\begin{equation}\label{eq:initial-gauge}
 \delta A=0,\qquad \delta B=BU,\qquad
 \delta C=[Z,C],\qquad \delta D=V,
\end{equation}
with $U,V,Z\in\slr$. The choice of $Z$ is allowed to change by an element
of $\hc$.

\smallskip\noindent\emph{The first band.}
Cyclic rotation of a product of invertible matrices preserves its
characteristic polynomial. Use the cyclic representative $A^kDC^lB$
for $u_0$. Its derivative is $A^kE_l$, where
\begin{equation}\label{eq:first-variation}
 E_l=(V+Z)C^lB+C^l(\Ad_B U-Z)B.
\end{equation}
For any $M$ and any $j$,
\begin{equation}\label{eq:Vandermonde}
 e_j(A^kM)=\sum_{|I|=j}a_I^k\det M[I,I].
\end{equation}
For fixed $j<r$ the distinct nonzero numbers $a_I$ make the coefficient
matrix for $k=1,\ldots,\binom rj$ a nonsingular Vandermonde matrix times
a nonsingular diagonal matrix. Thus stationarity for all $k\geq1$
forces every $j\times j$ principal minor of $M=C^lB$ to be stationary
under $E_l$. The full determinant is stationary too: the variations come
from products in $\SL_r$.

Positive powers of $C$ are Zariski dense in $\TC$. To see that positive
powers suffice, their Zariski closure in the connected torus is the same
as the closure of the full cyclic group; equivalently a Laurent
polynomial vanishing on all positive powers also vanishes on all integer
powers by the recurrence satisfied by its finitely many exponential
terms. Therefore the resulting regular identities hold with $C^l$
replaced by every $N\in\TC$.
Put
\[
 X=V+Z,\qquad Y=\Ad_BU-Z,\qquad W=X+Y=V+\Ad_BU.
\]
All principal minors of $NB$ are stationary for the variation
$(XN+NY)B$. At $N=I$, the one-by-one identities give $W_{ii}=0$.
Differentiate the stationary two-by-two identity along $N=I+tK+O(t^2)$,
where $K\in\hc$. The one-by-one stationary identities hold for every
$N$, so the derivative of the diagonal-product part is zero. The two
off-diagonal terms give, after dividing by $-B_{ii}B_{jj}\ne0$,
\begin{equation}\label{eq:pair-W}
 K_{ij}W_{ji}+K_{ji}W_{ij}=0.
\end{equation}
Condition (P1) forces $W_{ij}=W_{ji}=0$ for each pair. Hence $W=0$ and
$V=-\Ad_BU$, and the variation reduces to $[X,N]B$.

Choose a regular $N$ as in (P2). The matrix $NB$ has nonzero
off-diagonal entries, and each triangle difference is that of $N$
multiplied by the nonzero scalar $B_{ii}B_{jj}B_{kk}$.
Lemma~\ref{lem:minors} gives
\[
 [X,N]B=[h_0,NB]=[h_0,N]B
\]
for a diagonal $h_0$. Thus $X-h_0$ centralizes regular $N$, whose
centralizer in $\SL_r$ is $\TC$. Subtracting scalar parts makes
$h_0$ traceless and $X-h_0\in\hc$. Since $X=Z-\Ad_BU$, subtracting
the simultaneous diagonal conjugation by $h_0$ yields precisely
\begin{equation}\label{eq:residual-gauge}
 \delta A=0,\qquad \delta B=SB,\qquad
 \delta C=[S,C],\qquad \delta D=-S,
 \qquad S=\Ad_BU.
\end{equation}
Here diagonal conjugation does not change $A$ or $B$, and the remaining
part of $Z-h_0-S$ commutes with $C$.

The derivative of the surface relation in \eqref{eq:residual-gauge} is
\[
 (\Ad_A S-S)+(S-\Ad_C S)=(\Ad_A-\Ad_C)S.
\]
Consequently any relation tangent in this gauge satisfies
\begin{equation}\label{eq:T0}
 T_0:=\Ad_A S=\Ad_C S.
\end{equation}

\smallskip\noindent\emph{The second band.}
At the base point $u_1(k,l)$ has matrix $M=A^kBC^l$.
Differentiate the actual word $AB A^{k-1}CD C^{l-1}$ in the gauge
\eqref{eq:residual-gauge}. The $B$ contribution is $T_0M$.
For the last three factors one has
\[
 \delta(CDC^{l-1})=[S,C^l]-CSC^{l-1}.
\]
Using $CS=T_0C$ gives
\begin{align}
 \delta M
 &=T_0M+A^kB(S-T_0)C^l-MS\notag\\
 &=[T_0,M]+A^kB[R,C^l],\qquad R=S-\Ad_CS.\label{eq:second-variation}
\end{align}
The conjugation term has zero characteristic derivatives. Apply
\eqref{eq:Vandermonde} and density of the powers of $C$ once more.
All principal minors of $BN$ are stationary for the variation $B[R,N]$,
for every $N\in\TC$. The full determinant is stationary because
$\tr(N^{-1}[R,N])=0$. At a matrix from (P2),
Lemma~\ref{lem:minors} yields
\[
 B[R,N]=[h_1,BN]=B[h_1,N],
\]
so $R-h_1\in\hc$ after removing scalar parts. Hence
$(I-\Ad_C)S\in\h+\hc$. Combined with \eqref{eq:T0}, this is exactly
\eqref{eq:residual}. Proposition~\ref{prop:residual} forces $S=0$.
Undoing the gauges proves the kernel assertion. The converse holds because
all the functions are conjugation invariant.

\smallskip\noindent\emph{Smoothness and the quotient.}
Since $B$ has distinct diagonal entries, a $B$-invariant subspace is a
sum of coordinate lines. No proper nonzero such sum is invariant under
$C$, because every entry of $C$ is nonzero. The representation is
irreducible. The image of the derivative of the relation contains
\[
 \im(\Ad_A-I)+\im(\Ad_C-I),
\]
by left variations of $B$ and variations of $D$. Under the nondegenerate
trace pairing its annihilator is the common Lie centralizer
$\h\cap\hc=0$. The relation map is therefore submersive. The relation
variety is smooth of dimension $3(r^2-1)$ at the point; simultaneous
conjugation has dimension $r^2-1$. Its stabilizer is the scalar center,
which acts trivially. The standard irreducible smooth quotient description
from Section~\ref{subsec:geometry} gives a smooth character point of
dimension $2(r^2-1)$, with tangent space the quotient by conjugations.
The common-kernel assertion is thus equivalent to cotangent spanning.
\end{proof}

\begin{proposition}\label{prop:all-genus}
For every $g\geq2$, at the representation
\begin{equation}\label{eq:family}
 \rho_A=(A,B,C,I,I,\ldots,I)
\end{equation}
with the same open conditions, the functions in
Proposition~\ref{prop:core}, together with the traces of
$u_1(k,l)a_i$ and $u_1(k,l)b_i^{-1}$ for $i\geq3$, span the full
cotangent space of $\X$. A finite list of these functions suffices.
\end{proposition}
\begin{proof}
The commutator of an extra pair of identity matrices has zero first
derivative, so its two tangent variables $\delta A_i,\delta B_i\in\slr$
are unrestricted by the linearized relation. Infinitesimal simultaneous
conjugation is zero on these identity entries. Apply the core argument
to a common-kernel tangent and gauge its first four variations to zero.
The trace differentials of the additional curves then are exactly
\begin{equation}\label{eq:extra-pairings}
 \tr(U_{kl}\delta A_i),\qquad-\tr(U_{kl}\delta B_i),
 \qquad U_{kl}=A^kBC^l.
\end{equation}
No subtraction of a core trace is necessary: its derivative is already
zero in this gauge.

The matrices $U_{kl}$ for $1\leq k,l\leq r$ span all of $M_r(\C)$.
Indeed $A,A^2,\ldots,A^r$ span the diagonal matrices by Vandermonde
interpolation. For fixed row $i$,
\[
 \operatorname{row}_i(C^l)=
   \sum_{j=1}^r\mu_j^lQ_{ij}\operatorname{row}_j(Q^{-1}).
\]
All $Q_{ij}$ are nonzero and the $\mu_j$ are distinct, so varying
$l=1,\ldots,r$ spans all row vectors. The additional nonzero factor
$B_{ii}$ does not change the span. The trace pairing against all of
$M_r$ detects every tangent matrix, so \eqref{eq:extra-pairings} kills
all extra variations. Smoothness follows by the same relation
submersion and irreducibility argument as before. The quotient dimension
is $(2g-2)(r^2-1)$.

A spanning set in a finite-dimensional cotangent space has a finite
spanning subset. Choose one such subset at one point in the nonempty
open parameter set, and let $\F$ be the finite set of underlying simple
curves. Retain the required coefficient index for each chosen function;
we may add any of the core functions used in constructing this list.
No effective bound on this finite list is asserted.
\end{proof}
